\dykband{EXERCISES IN SWAFL}
{\small Open SWAFL (desktop or online) with the benchmark project. The chapter of the \href{https://github.com/halilunaziru73-creator/Coupled-Water-Fluxes-Pore-Scale-Dynamics-Precision-Irrigation-Mulching-Feedback-Loops/raw/main/Software/v1.1.0/SWAFL_Desktop/docs/manual_v1.1.0/SWAFL_Manual_v1.1.0.pdf}{SWAFL manual} that explains each task is given in brackets; record your values in a table and compare them with what is expected (\emph{Check}).}\par\vspace{3pt}
\begin{exlist}
\item {\color{khblue}[Manual, Chapter 21]} Open \emph{Crop water requirement} and choose maize, then wheat. Compare the $K_c$ curves and the seasonal water need.\par\noindent{\small\itshape Check: maize has a higher mid-season $K_c$ and a later peak.}
\item {\color{khblue}[Manual, Chapter 21]} In the same layer change the planting date by one month and record the change of the seasonal requirement.\par\noindent{\small\itshape Check: a later planting shifts the peak into the hotter period.}
\item {\color{khblue}[Manual, Chapter 19]} Import a weather file or use \emph{The weather of a place} for your own town and compare $ET_0$ with the benchmark record.\par\noindent{\small\itshape Check: the reference demand reflects your climate.}
\item {\color{khblue}[Manual, Chapter 22]} Open \emph{Reliability over 20 seasons} and read the years in which the crop was stressed. Relate them to the weather of those years.\par\noindent{\small\itshape Check: stress years coincide with dry, hot seasons.}
\end{exlist}
